\section{A finite-demand framework for first-order iPRA}
\label{sec:finite-demand-framework}

The central proposed interface is a theorem about a finite observation, accompanied
by an explicit list of the input observations it uses.  Infinite objects may occur
in the mathematical explanation, but the inductive argument is carried by a
primitive recursive state and a decidable invariant.  This is useful in two
ways.  It exposes exactly where an existence theorem needs a bound or a modulus,
and it prevents the considerable metatheory available in Lean from being
silently added to the object theory.

The transfer results below are proved directly in intuitionistic arithmetic.
Their role is narrower than a general conservation theorem: they compile a
specified class of finite arguments, with their hypotheses visible.  They do
not assert that every constructive proof, every computable procedure, or every
Lean theorem has a translation into \ipra.  The ordinary PR presentation and its
closure under bounded predicates are standard; see, for example,
\cite[Section~2]{Avigad2002}.  All the finite-demand arguments given here are
self-contained and do not invoke that paper's higher-type or classical
conservation results.

\subsection{The target theory and the permitted induction}

For this chapter, \ipra\ means first-order intuitionistic predicate logic with
equality, the natural-number constructors $0,S$, the successor axioms
\[
 Sx\ne0,\qquad Sx=Sy\longrightarrow x=y,
\]
and a function symbol for each fixed primitive recursive construction.  The
defining equations include composition and
\[
 r(0,\boldsymbol x)=g(\boldsymbol x),\qquad
 r(Sn,\boldsymbol x)=h(n,r(n,\boldsymbol x),\boldsymbol x).
\]
The induction schema is restricted to quantifier-free formulas $A$, with
parameters allowed:
\[
 \bigl(A(0,\boldsymbol x)\land
       \forall n(A(n,\boldsymbol x)\to A(Sn,\boldsymbol x))\bigr)
       \longrightarrow \forall n\,A(n,\boldsymbol x).
\]
Only finitely many of the PR symbols occur in any individual derivation.
A fresh name for a specified PR construction is legitimate in this all-PR
presentation.  This does not license totality of arbitrary terminating
programs.  When working in a smaller finite presentation, the introduction
must instead be justified by its established representation or translation
theorem.

There is no restriction saying that every theorem or every hypothesis must be
quantifier-free.  In particular, a statement such as
$\forall k\,C(a,k)$ is a legitimate hypothesis.  The restriction concerns the
formula to which induction is applied.  This distinction is indispensable:
one may instantiate a global regularity hypothesis at finitely many demanded
indices while using only a quantifier-free state invariant.

Integers, rationals, finite tuples, lists, and polynomials will always mean
natural-number codes.  Their operations are totalized on malformed codes; a
separate validity predicate specifies the intended inputs.  For example, a
rational is a pair $(z,d)$ with $z$ an integer code and $d>0$.  Equality is
\[
 (z,d)\simeq(w,e)\quad\Longleftrightarrow\quad ze=wd,
\]
so neither a quotient type nor a normalization algorithm is necessary.  Signs
and order can be tested by integer cross-multiplication with positive
denominators.  All these tests are primitive recursive.  An implementation
should keep validity predicates in its contracts even when its executable
functions return harmless default values on invalid data.

\subsection{Bounded truth is an arithmetic operation}
\label{sec:bounded-truth}

A \emph{bounded formula} in this chapter is built from equality of PR terms,
intuitionistic propositional connectives, and quantifiers of the form
$\forall u<t$ and $\exists u<t$, where $t$ does not contain the bound variable.
The bound may be any fixed PR term.  In particular, it need not be polynomial,
elementary, or practically computable.

\begin{lemma}[Booleanization of bounded formulas]
\label{lem:bounded-booleanization}
For every fixed bounded formula $B(\boldsymbol x)$ there is a PR function
$b_B(\boldsymbol x)\in\{0,1\}$ and an \ipra\ proof of
\[
 B(\boldsymbol x)\longleftrightarrow b_B(\boldsymbol x)=1.
\]
Consequently bounded formulas are decidable in \ipra, and induction on a
bounded formula is a derived rule of quantifier-free induction.
\end{lemma}

\begin{proof}
The elementary PR development of natural-number equality provides a Boolean
equality test.  Its usual correctness proof uses only the successor axioms and
quantifier-free induction; the same applies to truncated subtraction, order,
and the zero test.  Negation, conjunction, and disjunction are then Boolean
complement, product, and maximum.  For implication use the Boolean function
which returns $1$ when the antecedent test is $0$ and otherwise returns the
consequent test.  Its equivalence to intuitionistic implication is valid
because the antecedent is already decidable.  This step would be invalid for
an arbitrary, undecided proposition.

For bounded universal quantification use a prefix product, and for bounded
existential quantification use a prefix maximum.  Here is an explicit way to
justify the former without covert induction on a quantified formula.  Let
$c(i)\in\{0,1\}$ and define
\[
 p(0)=1,\qquad p(t+1)=p(t)c(t).
\]
Quantifier-free induction proves $p(t)\in\{0,1\}$ and $p(t+1)\le p(t)$.
With an arbitrary fixed $i$, another quantifier-free induction proves
\[
 i<t\land p(t)=1\quad\longrightarrow\quad c(i)=1.
\]
For the converse, fix $b$ and assume $c(i)=1$ for every $i<b$.  Induct on the
quantifier-free formula $t\le b\to p(t)=1$.  In the step with $t<b$, the fixed
hypothesis supplies $c(t)=1$.  Thus
\[
 p(b)=1\quad\longleftrightarrow\quad\forall i<b\;c(i)=1.
\]
The existential case can be proved by the explicit bounded search in
Lemma~\ref{lem:bounded-choice} below, or by the same argument with a running
maximum and a stored witness.  A structural induction on the fixed formula
$B$, in the metatheory, now gives its Boolean function and an \ipra\ proof of
its correctness.  Replacing $B(n)$ by $b_B(n)=1$ in the induction axiom proves
the last assertion.
\end{proof}

The adjective ``fixed'' matters.  The lemma supplies a schema indexed by a
finite formula and its fixed PR vocabulary.  It does not give a single PR truth
function for arbitrary formulas containing arbitrary PR program codes.  That
stronger assertion is false, as shown below.

\begin{lemma}[Bounded search and finite choice]
\label{lem:bounded-choice}
Let $R(\boldsymbol x,y)$ be a fixed PR relation.  There is a PR function
$\ell_R(\boldsymbol x,b)\le b$ such that \ipra\ proves
\begin{align*}
 \ell_R(\boldsymbol x,b)<b
   &\longleftrightarrow \exists y<b\;R(\boldsymbol x,y),\\
 \ell_R(\boldsymbol x,b)<b
   &\longrightarrow R(\boldsymbol x,\ell_R(\boldsymbol x,b)),\\
 y<\ell_R(\boldsymbol x,b)
   &\longrightarrow \neg R(\boldsymbol x,y).
\end{align*}
For a fixed PR bound $B(\boldsymbol x,i)$ and relation
$R(\boldsymbol x,i,y)$, a single PR finite-list function chooses all witnesses
below their bounds whenever
\[
 \forall i<n\;\exists y<B(\boldsymbol x,i)\;R(\boldsymbol x,i,y)
\]
holds.  No choice axiom is used.
\end{lemma}

\begin{proof}
Use $b$ as the sentinel value for failure and define
\[
 \ell(\boldsymbol x,0)=0,
\]
\[
 \ell(\boldsymbol x,b+1)=
 \begin{cases}
   \ell(\boldsymbol x,b),&\ell(\boldsymbol x,b)<b,\\
   b,&\ell(\boldsymbol x,b)=b\text{ and }R(\boldsymbol x,b),\\
   b+1,&\text{otherwise}.
 \end{cases}
\]
This is primitive recursion with decidable branches.  The three maintained
facts are that the stored value is at most the current bound; that a value
strictly below the bound is a successful witness; and that no earlier position
succeeds.  The last fact may be written as a prefix product of the Boolean
negations of $R$, so all three invariants are quantifier-free.

At the next stage, an already successful value is retained.  If none has yet
been found, the new position $b$ is tested.  A successful test stores $b$ and
uses the earlier failure product; an unsuccessful test extends that product.
This proves the invariants by quantifier-free induction.  In particular, the
sentinel occurs exactly when every position below $b$ fails.  Conversely, a
witness below $b$ precludes the sentinel: apply the earlier-failure invariant
to that witness.  This proves the displayed equivalence, constructively.

For finite choice, set
$w_i=\ell_{R(\boldsymbol x,i,-)}(B(\boldsymbol x,i))$ and build the finite
list $(w_0,\ldots,w_{n-1})$ by primitive recursion.  An entirely numerical
encoding is available.  Put
\[
 M=\max\bigl(1,\max_{i<n}B(\boldsymbol x,i)\bigr),\qquad
 W=\sum_{i<n}w_i(M+1)^i.
\]
Since $w_i\le M$, the base-$(M+1)$ digits of $W$ recover the $w_i$, and
$W<(M+1)^n$.  Digit extraction, the finite maximum, and the finite sum are all
PR.  Under the displayed hypothesis, each $w_i$ is strictly below its
corresponding bound and satisfies $R$.  The argument verifies a finite table;
it never chooses an unbounded sequence of witnesses.
\end{proof}

This lemma precisely identifies one frequent repair.  A finite family of
existence assertions causes no difficulty once PR bounds for the witnesses
have been proved.  Merely knowing $\forall i<n\,\exists y\,R(i,y)$ supplies no
such bound, and an unbounded search cannot be relabelled as bounded search.
Likewise, a terminating executable search has not yet been shown to be
primitive recursive: its termination argument must give an appropriate
bound or a PR construction of the output.

\subsection{The local demand theorem}
\label{sec:local-demand-theorem}

We next make the observation interface precise.  Fix the PR functions
\[
 \operatorname{init}(x),\quad
 \operatorname{step}(x,t,s),\quad
 k(x,t,s),\quad q(x,t,s,j),\quad \operatorname{ret}(x,N,s),
\]
and define the run by primitive recursion:
\[
 S(x,0)=\operatorname{init}(x),\qquad
 S(x,t+1)=\operatorname{step}(x,t,S(x,t)).
\]
The numerical parameter $x$ may contain all finite input data and parameters.
A query $q$ is simply a natural-number code for a demanded input fact.  Let
$U(x,q)$ be a fixed bounded predicate describing that fact.  At stage $t$,
the run demands $k(x,t,S(x,t))$ such facts.  Define
\begin{equation}
\label{eq:demand-log}
 L(x,N)\;:\!\Longleftrightarrow\;
 \forall t<N\;\forall j<k(x,t,S(x,t))\;
 U\bigl(x,q(x,t,S(x,t),j)\bigr).
\end{equation}
This is a bounded predicate, hence a single PR Boolean test.  Its query list
may be adaptive: later queries can depend on the already computed state.

\begin{theorem}[Finite-demand transfer]
\label{thm:finite-demand-transfer}
Let $P(x)$, $I(x,N,t,s)$, and $Q(x,N,z)$ be bounded predicates.  Suppose the
following three implications have \ipra\ derivations:
\begin{align}
 P(x)&\longrightarrow I(x,N,0,\operatorname{init}(x)),
   \label{eq:demand-base}\\
 t<N\land I(x,N,t,s)\land
 \bigl(\forall j<k(x,t,s)\;U(x,q(x,t,s,j))\bigr)
 &\longrightarrow
 I(x,N,t+1,\operatorname{step}(x,t,s)),
   \label{eq:demand-step}\\
 I(x,N,N,s)&\longrightarrow Q(x,N,\operatorname{ret}(x,N,s)).
   \label{eq:demand-finish}
\end{align}
Then \ipra\ proves
\begin{equation}
\label{eq:demand-conclusion}
 P(x)\land L(x,N)\quad\longrightarrow\quad
 Q\bigl(x,N,\operatorname{ret}(x,N,S(x,N))\bigr).
\end{equation}
The induction required by this proof is quantifier-free.
\end{theorem}

\begin{proof}
Fix $x,N$ and the two hypotheses in the antecedent of
\eqref{eq:demand-conclusion}.  Use the invariant
\[
 J(t)\;:\!\Longleftrightarrow\;
 t\le N\longrightarrow I(x,N,t,S(x,t)).
\]
By Lemma~\ref{lem:bounded-booleanization}, $J$ has a quantifier-free Boolean
form.  Its base follows from \eqref{eq:demand-base}.  To prove its step,
suppose $t+1\le N$.  Then $t<N$, so $J(t)$ gives the current invariant and
$L(x,N)$ gives exactly the finitely many hypotheses needed in
\eqref{eq:demand-step}.  Substitute the defining equation of $S$ to obtain
$J(t+1)$.  When $t+1>N$, the implication defining $J(t+1)$ is immediate.
Order is decidable, so this case distinction is intuitionistically valid.
Quantifier-free induction now gives $J(N)$; apply
\eqref{eq:demand-finish}.

The global statement $L(x,N)$ was a fixed hypothesis during the induction.
We did not induct on ``all inputs are valid'', on an unbounded universal
assertion, or on the correctness of arbitrary programs.  Every occurrence of
$U$ used by the proof is an entry of the finite demand log.
\end{proof}

\begin{corollary}[Global assumptions can be supplied locally]
\label{cor:global-to-local}
Under the hypotheses of Theorem~\ref{thm:finite-demand-transfer}, \ipra\ proves
\[
 P(x)\land\forall q\,U(x,q)
 \longrightarrow
 \forall N\;Q\bigl(x,N,\operatorname{ret}(x,N,S(x,N))\bigr).
\]
\end{corollary}

\begin{proof}
Fix $N$.  Instantiate the universal input assumption at each query in
\eqref{eq:demand-log}, obtaining $L(x,N)$, and apply the theorem.  Universal
generalization on $N$ completes the argument.  The quantifier-free induction
has already taken place inside the theorem.
\end{proof}

The corollary permits a convenient mathematical language involving globally
regular names.  The theorem is the more informative arithmetic statement:
it remains meaningful for a partial or inconsistent proposed object as long
as the particular finite observations required by this computation are
certified.  This is especially useful for interfaces between independently
developed theorem packets.

\subsection{Prefix contracts and their composition}

Many applications require all observations up to a stage rather than an
arbitrary query list.  Let $V(x,n)$ be a bounded prefix-validity predicate and
suppose \ipra\ proves
\[
 m\le n\land V(x,n)\longrightarrow V(x,m).
\]
An operation with an accuracy parameter $r$ has a \emph{prefix contract} when
there is a PR demand bound $\beta(x,r)$ and an \ipra\ proof of
\[
 V(x,\beta(x,r))\longrightarrow W(x,r,F(x,r)),
\]
where $F$ is PR and $W$ is bounded.  This contract states an actual finite
arithmetic implication.  Calling the output a real number, a function, a
root, a chain, or a geometric object adds no theorem until its intended
observations have been specified by $W$.

\begin{proposition}[Composition of prefix contracts]
\label{prop:prefix-composition}
Suppose an intermediate family of observations $A(x,j)$ satisfies
\[
 V(x,\beta(x,j))\longrightarrow A(x,j),
\]
and a second construction satisfies
\[
 \bigl(\forall j\le\gamma(x,r)\;A(x,j)\bigr)
       \longrightarrow B(x,r).
\]
All displayed matrices and functions are bounded or PR, and both implications
have \ipra\ proofs.  Then \ipra\ proves
\[
 V(x,\delta(x,r))\longrightarrow B(x,r),\qquad
 \delta(x,r)=\max_{j\le\gamma(x,r)}\beta(x,j).
\]
\end{proposition}

\begin{proof}
The finite maximum is PR.  For any $j\le\gamma(x,r)$, its defining property
gives $\beta(x,j)\le\delta(x,r)$.  Monotonicity of prefix validity and the
first contract imply $A(x,j)$.  Bounded universal introduction supplies the
premise of the second contract.  Its conclusion is $B(x,r)$.
\end{proof}

If $\beta(x,j)$ is monotone in $j$, the maximum can be replaced by
$\beta(x,\gamma(x,r))$.  Otherwise the maximum is necessary.  In particular,
a demand transformer must not be composed by blindly substituting precision
parameters when a later proof needs a whole prefix of earlier observations.
Adaptive dependencies are handled by Theorem~\ref{thm:finite-demand-transfer}
instead of by guessing a nonadaptive bound.

\subsection{What a proof compiler may legitimately promise}
\label{sec:proof-compiler}

A small source language can make the preceding discipline systematic.  Its
numeric computations are generated by fixed PR primitives, composition,
decidable conditionals, and primitive recursion on a numerical loop count.
Its assertions are bounded formulas.  Its proof rules are consequence,
sequential composition, decidable case analysis, and the loop rule with a
bounded invariant.  A primitive rule is available only together with an
\ipra\ derivation of its contract.  A mathematical premise such as regularity
of a named function must enter through a stated numerical observation
predicate, not as an unexamined side condition.

\begin{theorem}[Syntactic compilation of the finite proof discipline]
\label{thm:syntactic-compilation}
For every fixed well-formed source derivation built from the preceding rules,
there is an effective construction of an \ipra\ derivation of its concluding
contract.  Every numeric function introduced by the construction comes with
a finite PR definition.  Every inductive proof obligation is discharged by
quantifier-free induction after Booleanization.
\end{theorem}

\begin{proof}
Proceed by induction on the finite source derivation in the metatheory.
A primitive node already contains the required target derivation.  For
sequential composition, substitute the term denoting the first output in the
second contract and use intuitionistic implication elimination.  A
conditional is handled by decidability of its Boolean guard and disjunction
elimination; both branch functions and the combined conditional are PR.
A consequence node is ordinary intuitionistic implication composition.

For a loop, introduce a fresh function symbol with the displayed primitive
recursion equations for its state.  Booleanize the supplied bounded invariant
and apply Theorem~\ref{thm:finite-demand-transfer}.  The local premises used
by the body form the loop's finite demand log.  All list formation, finite
maximum, and Booleanization constructions use composition and primitive
recursion.  Thus the induction constructs an actual finite target derivation
at every node.

This is a construction on syntax.  The source derivation can be checked, its
new PR definition symbols can be emitted, and its target proof can be
assembled without evaluating all programs represented by source syntax.
Nothing in this construction supplies a universal semantic evaluator or a
uniform internal soundness principle for \ipra.
\end{proof}

The theorem is useful as a design specification for an eventual proof
compiler, but it is not a claim that such a compiler has been implemented in
this archive.  In particular, a collection of Lean theorem statements about
the intended semantic functions is not already an implementation of the
syntactic construction.  The output that ultimately matters is a derivation
in the chosen object-theory calculus, or a formally verified translation
into that calculus.

\subsection{Three distinct ways to represent a mathematical family}
\label{sec:three-representations}

The following distinctions remove a subtle source of overly strong claims.

\paragraph{A schema for a fixed PR name.}
Fix a particular PR function $a(k)$ and its definition.  A real-number name,
for example, can be described by rational approximations $a(k)$ and
arithmetical coherence assertions about pairs of indices.  An operation on
that name is another PR definition built from this particular $a$.  The
result is a theorem schema, one instance for every supplied PR definition.
There is no quantifier over all such functions inside first-order \ipra.

\paragraph{A uniformly PR numerical family.}
One may instead fix a jointly PR function $A(e,k)$, where $e$ is an ordinary
natural-number parameter.  Then \ipra\ can prove one theorem uniformly in
$e$.  This genuinely permits quantification over the parameter family, but
it covers only the family given by this one evaluator $A$.  It must not be
identified with an enumeration of all PR functions.

\paragraph{Certified finite observations.}
A third, more general, interface accepts a finite transcript.  For a fixed
register-machine or Turing-machine syntax, the relation
\[
 \operatorname{Trace}(e,u,v,t,c)
\]
asserts that $c$ is a length-at-most-$t$ computation of program $e$ on input
$u$ with output $v$.  Checking it is PR: parse the finite list of
configurations, verify the initial and final configurations, and verify each
adjacent transition by a bounded scan.  Equivalently, bounded execution with
an explicit numerical clock $t$ is PR.  A finite-demand theorem may
therefore be uniform in arbitrary program indices when it is conditional on
the relevant finite traces.

This last interface does not assert that every requested trace exists.  To
turn it into an unconditional theorem for a particular family, one must
supply a PR clock or a PR trace generator, and prove its correctness in the
object theory.  Passing the code of an arbitrary PR clock is insufficient:
the clock itself must be evaluated, which recreates the same issue.

\begin{proposition}[There is no universal PR evaluator]
\label{prop:no-universal-pr-evaluator}
There is no primitive recursive function $E(e,n)$ such that, for every unary
PR function $f$, some index $e$ satisfies
$E(e,n)=f(n)$ for all $n$.
\end{proposition}

\begin{proof}
If $E$ were PR, so would be $d(n)=E(n,n)+1$.  By the assumed universality,
there would be an index $e$ with $E(e,n)=d(n)$ for every $n$.  At $n=e$ this
gives $E(e,e)=E(e,e)+1$, a contradiction.
\end{proof}

This is a metamathematical diagonal argument about actual PR functions.  It
explains why even a perfectly decidable syntax for PR programs does not
justify a single PR evaluation operation for them all.  The distinction is
not a technical inconvenience to be hidden by notation; it is forced by the
mathematics.

Two useful operations nevertheless remain uniform.  A compiler can
transform finite program syntax into finite program syntax, and a checker
can verify a finite computation with a supplied clock.  Neither operation
evaluates arbitrary PR programs to completion.  The proposed framework uses
these two uniform operations together with fixed-name theorem schemas.

\subsection{Presentations, equality, and observations}
\label{sec:presentations}

A \emph{finite arithmetic presentation} consists of a PR validity predicate
$D(a)$, a PR equivalence test $a\simeq b$ on valid codes, and PR operation
codes with object-theory proofs of their required laws.  Rationals,
finite-dimensional rational vectors, finite polynomial expressions, and
finite combinatorial structures fit this format directly.  This does not
mean that every abstract countable structure has such a presentation; the
operations, validity tests, and equality laws are hypotheses to be supplied.

For completion-type objects, exact equality is usually an infinite relation.
A name $a$ can instead have observations $a[r]$ and finite consistency
predicates $V(a,n)$.  The informal assertion that it is a valid completed
object is
\[
 \operatorname{Valid}(a)\;:\!\Longleftrightarrow\;\forall n\;V(a,n).
\]
Here $a[r]$ means either a fixed PR name, a member of a uniformly PR family,
or a requested entry of a finite transcript, as distinguished above.
Arithmetic operations are supplied by finite contracts.  Equality of two
names is an observational assertion, such as an error inequality at every
accuracy; it is never assumed decidable merely because each individual
accuracy test is decidable.

For example, suppose a fixed PR family of rational approximations satisfies
\[
 \forall p,q\quad
 \abs{a[p]-a[q]}\le 2^{-p}+2^{-q}.
\]
For a finite demand $N$, all instances with $p,q\le N$ form a bounded
predicate after denominators are cleared.  A proof that a computed output
has a desired error at accuracy $r$ can request an appropriate finite
prefix.  Once that local contract is proved, the global coherence hypothesis
supplies it by instantiation.  There is no need to induct on the displayed
universal assertion or to construct a quotient of all Cauchy sequences
inside the arithmetic.

The same pattern handles many other mathematical concepts.  A continuous
map needs an approximation procedure and a supplied continuity modulus.
An invertible element needs an explicit inverse or an apartness bound
sufficient for the desired computation.  A compact set used for a search
needs a finite net with an explicit mesh and coverage contract.  An algebraic
extension can be presented by finite arithmetic data, but its field laws
still require a proof; a finite multiplication table over an infinite base
field is not automatically a field certificate.  These are interfaces with
mathematical content, not declarations that erase the original hypotheses.

\subsection{A reusable algebraic application: polynomial certificates}
\label{sec:polynomial-certificates}

The restriction against a universal PR evaluator does \emph{not} prohibit
uniform evaluation of a fixed expression grammar.  This distinction gives a
particularly useful reusable module for algebraic and geometric packets.

Fix a PR-presented commutative ring: its carrier predicate $D$, equality
$\simeq$, and operations $0,1,+,-,\cdot$ are fixed PR data, and their carrier-closure, ring,
and congruence laws have \ipra\ proofs.  Expressions are finite trees built
from variables, integer constants, addition, negation, and multiplication.
Let $\operatorname{Eval}(e,a)$ evaluate the tree using a finite assignment
$a$ of carrier elements.  Let $\operatorname{NF}(e)$ be the usual sparse
commutative polynomial obtained by expanding products, representing each
monomial by a finite exponent vector (padding with zeros when necessary),
ordering these vectors, and collecting equal monomials.  Integer coefficients are stored as integer codes.

\begin{proposition}[Uniform soundness of a fixed algebraic certificate language]
\label{prop:polynomial-normalization}
There are PR functions $\operatorname{Eval}$ and $\operatorname{NF}$, uniform
in expression codes and finite assignments, for which \ipra\ proves
\[
 \operatorname{Good}(e,a)\quad\longrightarrow\quad
 \operatorname{Eval}(e,a)\simeq
 \operatorname{EvalPoly}(\operatorname{NF}(e),a).
\]
Consequently equality of normal forms is a sound PR certificate for equality
of the two expressions, on valid assignments.
\end{proposition}

\begin{proof}
Choose a tree coding in which every immediate subtree has smaller numerical
code than its parent.  Such a coding is obtained, for example, by positive
tagging of a pairing function that dominates its arguments.  Parsing and
checking this property are bounded operations.  Evaluation can now be
computed by course-of-values recursion on the numerical code: keep a finite
table of values for all smaller codes, and use the already stored child
values at the current code.  Invalid codes receive the default value $0$.
The table step uses only the fixed ring operations, so the table and its
lookup operation are PR.  This may be an inefficient evaluation method,
but its PR status is immediate.

Normal forms are computed by the same table construction.  A constant is a
one-term polynomial, and a variable is a one-term monomial.  Addition
concatenates the two coefficient lists and collects equal exponent vectors.
Negation changes coefficient signs.  Multiplication forms the finite list
of pairwise products: coefficients multiply and exponent vectors add.
Insertion sorting and collection are finite recursions, hence PR.  The
lengths and coefficient sizes can grow very rapidly, but all the stated
operations are finite compositions and primitive recursions.

We first verify the finite-list identities used by these operations.
Associativity proves that concatenation evaluates as a sum.  Commutativity
proves that adjacent swaps preserve the sum, so insertion sorting preserves
it.  Distributivity and repeated addition prove that merging equal monomials
adds their integer coefficients without changing evaluation.  For products,
induction on the first list and then on the second list gives
\[
 \left(\sum_{i<m}u_i\right)\left(\sum_{j<n}v_j\right)
       =\sum_{i<m}\sum_{j<n}u_i v_j.
\]
The induction formulas compare two coded ring elements, so they are
quantifier-free.  The identity $a^{r+s}=a^r a^s$ is obtained by
quantifier-free induction on $s$; applying it to the finitely many variables
shows that adding exponent vectors multiplies monomial evaluations.
Negative coefficients are handled by the additive inverse laws.  Each
identity is interpreted using $\simeq$, with the supplied congruence laws
justifying substitution.

Finally, prove the normalization claim simultaneously for all valid
expression codes below a numerical cutoff $b$.  The invariant is
\[
 \forall e<b\quad
 \bigl(\operatorname{Good}(e,a)\to
 \operatorname{Eval}(e,a)\simeq
 \operatorname{EvalPoly}(\operatorname{NF}(e),a)\bigr).
\]
It is a bounded formula and therefore has a quantifier-free Boolean form.
At a new code, every child code is smaller, so the induction hypothesis
applies to all children.  The finite-list identities just proved establish
the claim for the corresponding constructor.  An invalid code imposes no
obligation.  This proves the invariant by permitted induction, and hence the
claim for every expression code.  If two normal forms are equal as finite
codes, substitute that equality and use symmetry and transitivity of
$\simeq$ to obtain equality of the original expressions.
\end{proof}

This is an internal soundness theorem for one restricted certificate
language, justified by a specific PR evaluation and a specific bounded
invariant.  It is not uniform reflection for arithmetic proofs.  The
expression grammar has a fixed finite set of already defined operations;
it contains no constructor that executes an arbitrary coded primitive
recursion program.

The proposition is a common foundation for exact trigonometric polynomial
identities, rational-coordinate geometry identities, determinant identities
at a fixed dimension, and many of the repository's certificate-based
arguments.  It proves \emph{soundness} of certificates.  A uniform theorem
about all dimensions or all degrees additionally needs a PR certificate
generator and an \ipra\ proof that its outputs pass the checker, or a direct
uniform algebraic argument.  Running a generator on finitely many examples
cannot establish that extra statement.

\subsection{The Lean boundary and the object-theory goal}
\label{sec:lean-boundary}

A complete formalization should separate four layers.

\begin{enumerate}
\item \textbf{Executable finite arithmetic.}  Implement natural-number
codes, rational and integer operations, finite lists, Boolean predicates,
state recursions, and demand transformers.  These are ordinary Lean
functions at this stage.
\item \textbf{Mathematical correctness.}  Prove the local numerical
contracts and the finite-demand lemmas in Lean.  This validates the
mathematical arguments, but Lean's unrestricted induction is not itself an
\ipra\ derivation.
\item \textbf{PR syntax and bounded formulas.}  Give every used executable
function an explicit construction in a syntactic PR language; relate that
construction to the executable function.  Define object-language terms,
formulas, substitution, the quantifier-free test, and derivations.  Compile
bounded predicates to Boolean PR terms with derivations of the
corresponding equivalences.
\item \textbf{Object-theory derivations.}  Construct the syntactic proof of
the relevant arithmetic contract.  If the final target is ConciseIPRA,
compose with an established translation into its native derivability
relation, with all hypotheses of that translation supplied.
\end{enumerate}

The principal compiler theorem has the schematic target
\[
 \operatorname{SourceDerivation}(d,\Gamma,\varphi)
 \longrightarrow
 \operatorname{Derivation}_{\ipra}
       (\operatorname{compile}(d),\Gamma^*,\varphi^*),
\]
where the stars denote an explicitly defined syntax translation.  Its
conclusion is a proposition \emph{in Lean about finite derivation objects}.
The step from that conclusion to a native ConciseIPRA derivation must use a
proved bridge, not an assumed semantic reflection rule.

For the finite-demand theorem, the first genuinely arithmetic proof object
should use precisely the Booleanization of
\[
 t\le N\to I(x,N,t,S(x,t))
\]
as its induction formula.  Its base and step are obtained by substitution
from the supplied base and step derivations; its final inference substitutes
$N$ and uses the finishing derivation.  These instructions fix the intended
object-level proof rather than merely asking a theorem prover to rediscover
it.

For bounded search, the proof object should carry the running sentinel,
its success test, and the prefix-failure product.  For prefix composition,
it should carry the PR finite maximum and its elementary upper-bound
lemma.  For polynomial normalization, it should carry the bounded-prefix
invariant over subexpression codes displayed above.  These are concrete
main lemmas with finite syntax, explicit induction formulas, and clear
implementation paths.

Finally, the assertion that a proof checker is PR does not justify
\[
 \operatorname{Check}(p,\ulcorner\varphi\urcorner)=1\longrightarrow\varphi
\]
uniformly for all arithmetic formulas inside the theory being checked.
Such a statement is an additional reflection principle; at falsity it
already supplies the theory's own consistency assertion.  Restricted
reflection for a particular finite algebraic language can be established by
the preceding normalization proof.  General arithmetic reflection cannot
be silently substituted for that proof.  This boundary is essential when
reusing repository results whose Lean statements already live in a rich
ambient metatheory.
